% Ilustraciones independientes: aquí si caben; de lo contrario, arriba o
% abajo de una página cercana. Los capítulos vacían las figuras pendientes;
% las secciones dejan continuar el texto para no crear huecos al vaciarlas.
% Los diagramas de las demostraciones permanecen en el flujo matemático.
\setcounter{topnumber}{2}
\setcounter{bottomnumber}{2}
\setcounter{totalnumber}{3}
\renewcommand{\topfraction}{.9}
\renewcommand{\bottomfraction}{.85}
\renewcommand{\textfraction}{.1}
\renewcommand{\floatpagefraction}{.7}
\setlength{\textfloatsep}{14pt plus 3pt minus 2pt}
\setlength{\intextsep}{12pt plus 2pt minus 2pt}
\setlength{\floatsep}{12pt plus 2pt minus 2pt}
\captionsetup[figure]{font=small,labelfont={bf,color=rojoCIMAT},
  textfont={color=grisCIMAT!80!black},labelsep=period,skip=7pt}
\crefname{figure}{figura}{figuras}
\Crefname{figure}{Figura}{Figuras}

% Uso: \figuraApuntes{Figuras/archivo}{fig:etiqueta}{Título breve}
\newcommand{\figuraApuntes}[3]{%
  \par
  \begin{figure}[!htbp]
    \centering
    \input{#1}%
    \caption{#3}\label{#2}%
  \end{figure}%
  \par
}
